% answers_short.tex —— 1986 年·广东卷·文史类 · 简短答案速查
% 由 tools/make_short_answers.py 自 answers.tex 的 \gksolhead/\gkanswer 提取，
% 并与 1986/1986_数学真题_广东文_答案要点.md 逐题比对后生成，勿手改（改源 MD 后重跑）。

\begin{gknotice}
\textbf{使用说明}：本节只给最终答案便于快速核对，完整推导见解析版。答案为编写者独立推导，\textbf{非官方标准答案}。
\end{gknotice}

\gkpart{大题}
\begin{itemize}[leftmargin=2.6em]
  \item[\textbf{1.}] (1) D；(2) C；(3) E；(4) D；(5) A；(6) B；(7) A；(8) D；(9) E；(10) A；(11) E；(12) C；(13) B；(14) E；(15) B；(16) A；(17) B；(18) D；(19) B；(20) D；(21) B；(22) C；(23) E；(24) C；(25) E（固定源图形选项缺失，独立图形为左下四分之一弧）
  \item[\textbf{2.}] (1) $y=(\sqrt3)^x$；(2) $x^2=y$（固定源乘号为缺字方框）；(3) $\dfrac32$；(4) $4$；(5) $[-2,-1)$；(6) $180$；(7) $243$；(8) $\dfrac12$；(9) 待核：按固定源题面不存在实等差数列；解答页列 $180$；(10) $-2$
  \item[\textbf{3.}] (1) $x=2+\sqrt3$；(2) $a>0$，且唯一解为 $\displaystyle x=\frac{a+2+\sqrt{a^2+4a}}2$
  \item[\textbf{4.}] $\theta=\dfrac{3\pi}4$ 时最大值为 $3$；$\theta=\dfrac\pi4$ 时最小值为 $\dfrac13$
  \item[\textbf{5.}] 见解析正文。
  \item[\textbf{6.}] (1) $\displaystyle \frac{\sqrt{20b^2-5a^2}}{5a}$；(2) $\displaystyle \frac a{16}\sqrt{8b^2-2a^2}$
  \item[\textbf{7.}] $\displaystyle -\frac{2\sqrt{13}}{13}<m<\frac{2\sqrt{13}}{13}$
\end{itemize}
